% !TeX program = lualatex
% Definition and sources only; this file is not an LLM solution attempt.
\documentclass[11pt,a4paper]{article}
\usepackage[margin=25mm]{geometry}
\usepackage{fontspec}
\setmainfont{lmroman10-regular.otf}[BoldFont=lmroman10-bold.otf,
 ItalicFont=lmroman10-italic.otf,BoldItalicFont=lmroman10-bolditalic.otf]
\usepackage{amsmath,amssymb,mathtools}
\usepackage{unicode-math}
\setmathfont{latinmodern-math.otf}
\AtBeginDocument{\let\setminus\smallsetminus}
\usepackage{xurl,hyperref}
\hypersetup{colorlinks=true,urlcolor=blue}
\setcounter{secnumdepth}{0}
\setlength{\parindent}{0pt}
\setlength{\parskip}{5pt}
\setlength{\emergencystretch}{3em}
\begin{document}
\section*{Questions in Geometric Group Theory — Q 2.11}\label{statement}

\subsection{Definitions and mathematical statement}
A proper metric space has compact closed bounded balls. A geodesic metric space is CAT$(0)$ if its geodesic triangles satisfy Euclidean comparison: distances between side points are no greater than the corresponding comparison distances. Let a finitely generated discrete group $G$ act on a proper CAT$(0)$ space $X$ by isometries. Require proper discontinuity in the compact-set sense:
\[
 \{g\in G:gK\cap K\ne\varnothing\}\text{ is finite}
 \quad\text{for every compact }K\subset X.
\]
In particular, every point stabilizer is finite.

(a) Can $G$ be an infinite torsion group, meaning that every element has finite order although the group itself is infinite?

(b) If the action is also cocompact, meaning $X/G$ is compact, can $G$ contain an infinite torsion subgroup $H$? The subgroup in (b) is not assumed finitely generated, and its restricted action need not be cocompact.

Both questions permit arbitrarily large finite orders; torsion does not mean that a single exponent annihilates every element. The space may have infinite dimension, but it remains proper. A finite torsion group poses no obstruction and is not the requested example. Conversely, constructing an infinite torsion action with an infinite stabilizer, or on a nonproper space, would fail the stated hypotheses. The issue in the second part is the subgroup structure of a group acting cocompactly, rather than whether that ambient group itself has only torsion elements.
\subsection{Short English statement}
Can infinite torsion occur in a proper CAT(0) action, either as the whole finitely generated acting group or as a subgroup under cocompactness?
\subsection{Sources}
\begin{itemize}
\item[S1] ULAM.AI and the UnsolvedMath Contributors, supplied problems.json, catalogue entry, AMR Open Problem Lists. Catalogue identity and source transcription; CC BY 4.0.
\url{https://huggingface.co/datasets/ulamai/UnsolvedMath}.
\item[S2] Bestvina - Questions in Geometric Group Theory (pdf) (2004), Question 2.11 (PDF page 8). 
\url{https://www.math.utah.edu/~bestvina/eprints/questions-updated.pdf}.
\end{itemize}
\end{document}
